% from ega3-1-fr.tex
\oldpage[0\textsubscript{III}]{22}
résultera que $\widehat A'$ est $A'$-plat
(\hyperref[0.6.2.3-fr]{$0_{\mathrm I}$, 6.2.3}), et comme
$\mathfrak m'\widehat A'\ne\widehat A'$, $\widehat A'$ est un $A'$-module
fidèlement plat (\hyperref[0.6.6.2-fr]{$0_{\mathrm I}$, 6.6.2}), d'où on
conclura finalement que $A'$ est \emph{noethérien}
(\hyperref[0.6.5.2-fr]{$0_{\mathrm I}$, 6.5.2}), ce qui achèvera la
démonstration du lemme.

On a $\widehat A'=\varprojlim_n A'/\mathfrak m'^{n}$~; en raison de ce que
$A'$ est $A_\lambda$-plat, on a
\[
  \mathfrak m'^{n}/\mathfrak m'^{n+1}
  =(\mathfrak m_\lambda^n/\mathfrak m_\lambda^{n+1})
    \otimes_{A_\lambda}A'
  =(\mathfrak m_\lambda^n/\mathfrak m_\lambda^{n+1})
    \otimes_{K_\lambda}(K_\lambda\otimes_{A_\lambda}A')
  =(\mathfrak m_\lambda^n/\mathfrak m_\lambda^{n+1})
    \otimes_{K_\lambda}K~;
\]
comme $\mathfrak m_\lambda^n/\mathfrak m_\lambda^{n+1}$ est un
$K_\lambda$-espace vectoriel de dimension finie,
$\mathfrak m'^{n}/\mathfrak m'^{n+1}$ est un $K$-espace vectoriel de
dimension finie pour tout $n\geq0$. Il résulte donc de
(\hyperref[0.7.2.12-fr]{$0_{\mathrm I}$, 7.2.12}) et
(\hyperref[0.7.2.8-fr]{$0_{\mathrm I}$, 7.2.8}) que $\widehat A'$ est
noethérien. On sait en outre que l'idéal maximal de $\widehat A'$ est
$\mathfrak m'\widehat A'$ et que
$\widehat A'/\mathfrak m'^{n}\widehat A'$ est isomorphe à
$A'/\mathfrak m'^{n}$~; comme
$A'/\mathfrak m'^{n}=(A_\lambda/\mathfrak m_\lambda^n)
\otimes_{A_\lambda}A'$, $A'/\mathfrak m'^{n}$ est un
$(A_\lambda/\mathfrak m_\lambda^n)$-module plat
(\hyperref[0.6.2.1-fr]{$0_{\mathrm I}$, 6.2.1})~; le critère
(\hyperref[0.10.2.2-fr]{10.2.2}) est donc applicable à la
$A_\lambda$-algèbre noethérienne $\widehat A'$, et montre que $\widehat A'$
est $A_\lambda$-plat. C.Q.F.D.

\begin{env}[10.3.1.4]
\label{0.10.3.1.4-fr}
Abordons maintenant le cas général. Il existe un ordinal $\gamma$ et pour
tout ordinal $\lambda\leq\gamma$ un sous-corps $k_\lambda$ de $K$ contenant
$k$, tels que~: $1^\circ$ Pour tout $\lambda<\gamma$, $k_{\lambda+1}$ soit
une extension de $k_\lambda$ engendrée par un seul élément~; $2^\circ$ Pour
tout ordinal $\mu$ sans prédécesseur,
$k_\mu=\bigcup_{\lambda<\mu}k_\lambda$~; $3^\circ$ $K=k_\gamma$. Il suffit en
effet de considérer une bijection $\xi\mapsto t_\xi$ de l'ensemble des
ordinaux $\xi\leq\beta$ (pour un $\beta$ convenable) sur $K$, de définir
$k_\lambda$ par induction transfinie (pour $\lambda\leq\beta$) comme la
réunion des $k_\mu$ pour $\mu<\lambda$ si $\lambda$ n'a pas de prédécesseur,
et, si $\lambda=\nu+1$, comme $k_\nu(t_\xi)$, où $\xi$ est le plus petit
ordinal tel que $t_\xi\notin k_\nu$~; $\gamma$ est alors par définition le plus
petit ordinal $\leq\beta$ tel que $k_\gamma=K$.

Cela étant, nous allons définir, par récurrence transfinie, une famille
d'anneaux locaux noethériens $A_\lambda$ pour $\lambda\leq\gamma$, et des
homomorphismes locaux $f_{\mu\lambda}:A_\lambda\to A_\mu$ pour
$\lambda\leq\mu$, vérifiant les conditions suivantes~:
\begin{enumerate}
  \item[(i)] $(A_\lambda,f_{\mu\lambda})$ est un système inductif et $A_0=A$.
  \item[(ii)] Pour tout $\lambda$, on a un $k$-isomorphisme
    $A_\lambda/\mathfrak J A_\lambda\xrightarrow{\sim}k_\lambda$.
  \item[(iii)] Pour $\lambda\leq\mu$, $A_\mu$ est un $A_\lambda$-module plat.
\end{enumerate}

Supposons donc les $A_\lambda$ et les $f_{\mu\lambda}$ définis pour
$\lambda<\mu<\xi$, et supposons en premier lieu que $\xi=\zeta+1$, de sorte
que $k_\xi=k_\zeta(t)$. Si $t$ est transcendant sur $k_\zeta$, on définit
$A_\xi$ suivant le procédé de (\hyperref[0.10.3.1.1-fr]{10.3.1.1}) comme
égal à $(A_\zeta[t])_{\mathfrak J A_\zeta[t]}$~;
$f_{\xi\zeta}$ est l'application canonique, et pour $\lambda<\zeta$, on prend
$f_{\xi\lambda}=f_{\xi\zeta}\circ f_{\zeta\lambda}$~; la vérification des
conditions (i) à (iii) est alors immédiate, vu ce qui a été démontré en
(\hyperref[0.10.3.1.1-fr]{10.3.1.1}). Supposons ensuite $t$ algébrique, et
soient $h$ son polynôme minimal dans $k_\zeta[T]$, $H$ un polynôme unitaire de
$A_\zeta[T]$ dont l'image dans $k_\zeta[T]$ est $h$~; on prend alors $A_\xi$
égal à $A_\zeta[T]/(H)$, les $f_{\xi\lambda}$ se définissant comme
précédemment~; la vérification des conditions (i) à (iii) résulte alors de
ce qui a été vu en (\hyperref[0.10.3.1.2-fr]{10.3.1.2}).

Supposons maintenant que $\xi$ n'ait pas de prédécesseur~; on prend alors
pour $A_\xi$ la limite inductive du système inductif d'anneaux locaux
$(A_\lambda,f_{\mu\lambda})$ pour $\lambda<\xi$~; $f_{\xi\lambda}$ est définie
comme l'application canonique pour $\lambda<\xi$. Le fait que $A_\xi$ soit
local noethérien, que les $f_{\xi\lambda}$ soient des homomorphismes locaux,
et les conditions (i) à (iii) pour $\lambda\leq\xi$
